\label{chap:83-14}

\section{Derivación de los canales \(90/120\) mediante la inclusión de incidencias hexada--octada}

Fijados la dodecada y el alineamiento del
\cref{p12-83-c17-s1:sec:w12-w24-elevacion}, sea \(H\) una hexada en las coordenadas HMT,
sea \(\ell(H)\subset D\) su imagen y sea
\(O_H=\iota_D(\ell(H))\) su elevación única. Al marcar un punto
y un par de puntos de la hexada residual se obtienen
\[
F_6(H)=\ell(H)\times\binom{\ell(H)}2,
\qquad
|F_6|=6\binom62=90.
\]
Si el punto marcado puede recorrer la octada:
\[
F_8(H)=O_H\times\binom{\ell(H)}2,
\qquad
|F_8|=8\binom62=120.
\]
La inclusión \(\ell(H)\hookrightarrow O_H\) induce el morfismo de incidencia
\[
 \mathcal I_{6\to8}:F_6(H)\hookrightarrow F_8(H).
\]
Al normalizar la diferencia por el factor
elegido \(24\binom62\), se obtiene
\[
\boxed{
\frac{90-120}{24\binom62}
=\frac{6-8}{24}
=-\frac1{12}.}
\]
La igualdad es un defecto normalizado de la interfaz combinatoria
\(6\to8\). El factor \(24\binom62\) forma parte de la definición de esta
normalización; la inclusión por sí sola determina la diferencia \(-30\).
El operador \(\mathcal I_{6\to8}\) tiene como dominio configuraciones
marcadas de incidencia; no es el extractor transversal del estado
dodecafásico ni la familia de momentos angulares.

\section{Transporte de \(-1/12\) a los canales \(90/120\) en \(5\) realizaciones algebraicas}

Además del defecto normalizado de la inclusión hexada--octada, aparecen
cinco realizaciones operatoriales del mismo racional. Sus dominios se
especifican por separado; la igualdad de valores no identifica los
operadores sin un entrelazador adicional.

\subsection{Diferencia de periodicidades dodecafásicas}

Un diente de \(\Cdoce\) mide \(30^\circ\). Entonces
\[
\frac{30}{120}-\frac{30}{90}
=\frac14-\frac13
=-\frac1{12}.
\]
Es la diferencia orientada entre los periodos de \(90^\circ\) y
\(120^\circ\) en esta realización.

\subsection{Pseudoinversa del operador de Gram sobre la componente viva}

En el salto exacto \(60\to81\), dos fibras tienen tamaños \(0\) y
\(12\). El operador de Gram es
\[
K_{60,81}=\operatorname{diag}(0,12).
\]
Sobre la componente no nula de su rango, la pseudoinversa vale \(1/12\); con la orientación
negativa de borde,
\[
E_{\rm surv}=-K_{60,81}^{+}
\]
produce \(-1/12\).

\subsection{2.ª diferencia entre escalas triádicas}

Definamos
\[
T_L=\sum_{n=1}^{L-1}n(L-n)=\frac{L(L^2-1)}6,
\qquad
a(L)=-\frac{T_L}{L^2}=-\frac L6+\frac1{6L}.
\]
La primera resta de escala es
\[
A_1(L)=a(L)-\frac{a(3L)}3=\frac4{27L},
\]
y la segunda
\[
C(L)=LA_1(L)-\frac{3L}{3}A_1(3L)=\frac8{81}.
\]
El residuo \(8/81\) es independiente de \(L\). Para obtener
\(-1/12\) se aplica la normalización declarada
\[
R(L)=-\frac{27}{32}C(L)=-\frac1{12}.
\]
El cálculo deriva \(8/81\); no deriva por sí solo el factor \(-27/32\).

\subsection{Composición de transformaciones modulares y preservación del estado firmado}

La función eta satisface
\[
\frac{\eta(\tau)}{\eta(3\tau)}
=q^{-1/12}\prod_{n\ge1}(1+q^n+q^{2n})^{-1}.
\]
El exponente procede de \(1/24-3/24=-1/12\). Al tomar doce copias:
\[
t_3(q)=\left(\frac{\eta(\tau)}{\eta(3\tau)}\right)^{12}
=q^{-1}-12+54q-76q^2-243q^3+\cdots.
\]
El coeficiente \(54\) coincide con el defecto APP\@. La identidad racional
\[
j=\frac{(t_3+27)(t_3+243)^3}{t_3^3}
\]
establece la conexión con la rama modular clásica.

\subsection{Defecto de los proyectores de diferencia cíclica}

Sea \(S\) el desplazamiento cíclico sobre \(\QQ^{12}\) y definamos
\[
 D_r=I-S^r.
\]
El núcleo de \(D_r\) está formado por los vectores constantes en cada órbita
de \(S^r\); por tanto,
\[
 \dim\ker D_r=\gcd(12,r).
\]
Sean \(\Pi_{\ker D_3}\) y \(\Pi_{\ker D_4}\) los proyectores racionales sobre
\(\ker D_3\) y \(\ker D_4\), y sea
\(\tau(T)=\operatorname{Tr}(T)/12\) la traza normalizada. Entonces
\[
 \boxed{
 \tau(\Pi_{\ker D_3}-\Pi_{\ker D_4})
 =\frac{\operatorname{rank}\Pi_{\ker D_3}
       -\operatorname{rank}\Pi_{\ker D_4}}{12}
 =\frac{3-4}{12}
 =-\frac1{12}.}
\]
Esta realización expresa el racional como defecto transversal de los pasos
tres y cuatro en el ciclo de doce fases.

\begin{remark}[Dominios de las realizaciones]
Son exactos el defecto APP \(54\), el exponente eta \(-1/12\), el defecto
combinatorio de la elevación hexada--octada, y las
identidades modulares una vez fijado \(t_3\). También son exactos el valor
de la pseudoinversa, el extractor normalizado y el defecto de proyectores.
Cada igualdad conserva su dominio y su aplicación correspondiente.
\end{remark}


\section{Parámetro de Barbero--Immirzi: operador de área, espectro,
entrelazador de canales y entropía modular}
\label{p12-83-c14-s2:sec:area-entrelazamiento-barbero-rev11}
\index{Barbero--Immirzi!operador de área}
\index{entropía modular!Hamiltoniano modular}

La identificación HMT del funcional \(\Gamma_{6\to8}\) con la coordenada
de Barbero--Immirzi se realiza sobre el mismo borde triádico que soporta la
incidencia (90/120). El bulk y su hoja de borde poseen igual cardinalidad,
\[
 \mathcal B=(\mathbb Z/9\mathbb Z)^3,
 \quad |\mathcal B|=729,
 \qquad
 \partial\mathcal B=(\mathbb Z/27\mathbb Z)^2,
 \quad |\partial\mathcal B|=729,
\]
mediante el reagrupamiento de seis trits en tres coordenadas nonádicas o dos
coordenadas de orden veintisiete.

\begin{theorem}[Ley areal del corte triádico]
\label{p12-83-c14-s2:thm:ley-areal-corte-triadico-rev11}
En el grafo cúbico \(N\times N\times N\) con capacidad unitaria, el corte
mínimo entre dos caras opuestas tiene capacidad \(N^2\). Para \(N=9\),
\[
 \mathcal A_{\min}=81,
 \qquad S_{\rm tri}=\mathcal A_{\min}\log3=81\log3.
\]
\end{theorem}
\begin{proof}
Las \(N^2\) rectas paralelas a la dirección normal son caminos disjuntos en
aristas, por lo que todo corte las intercepta. Una capa transversal contiene
exactamente \(N^2\) aristas y alcanza la cota. Cada enlace triádico aporta
\(\log3\) a la entropía de un estado máximamente mezclado.
\end{proof}

En el toro de borde \(27\times27\), sea \(A\) el bloque canónico
\(9\times9\). Su frontera se descompone de manera exacta en
\begin{equation}
 \partial A=\partial A_{90}\sqcup\partial A_{120},
 \qquad |\partial A_{90}|=|\partial A_{120}|=18.
 \label{p12-83-c14-s2:eq:corte-90-120-rev11}
\end{equation}
Los enlaces horizontales representan el canal (90), y los verticales, el
canal (120). Esta descomposición convierte la pareja de incidencias en una
descomposición operatoria del borde.

La entropía areal \(S_{\rm tri}=81\log3\) pertenece al estado máximamente
mezclado de los enlaces del corte. La entropía bicapa
\(S_{\rm bi}(y)\) de la \cref{def:entropia-bicapa-rev6} pertenece a la
distribución normalizada entre las dos hojas; ambos objetos conservan
dominios distintos.

Sea \(N=\operatorname{diag}(0,1,2)\) el operador número de un enlace trítico.
La realización areal declara la escala positiva
\[
 \theta_{\rm clk}=\frac{C^*}{6}\frac{\pi}{180},
 \qquad
 a_0=\frac{1+2\cos(10^\circ)}{3}\,\theta_{\rm clk},
 \qquad \gamma_{\rm HMT}=\Gamma_{6\to8}.
\]
La fórmula de \(a_0\) constituye la estructura de partida de esta
realización; fijada esa escala, el operador y su espectro se deducen
exactamente.

\begin{definition}[Operador de área HMT]
Sobre la hoja de borde
\(\mathcal H_{\partial}=\bigotimes_{e\in\partial A}\mathbb C^3\), se define
\begin{equation}
 \widehat{\mathcal A}_{\rm HMT}
 =\gamma_{\rm HMT}a_0\sum_{e\in\partial A}N_e.
 \label{p12-83-c14-s2:eq:operador-area-hmt-rev11}
\end{equation}
Su espectro en el corte canónico es
\[
 \operatorname{Spec}\widehat{\mathcal A}_{\rm HMT}
 =\{n\gamma_{\rm HMT}a_0:0\leq n\leq72\},
\]
con multiplicidades dadas por los coeficientes de
\((1+z+z^2)^{36}\). La prolongación por cortes compatibles produce la
familia protegida \(\mathcal A_n^{\rm prot}=n\gamma_{\rm HMT}a_0\).
\end{definition}

El entrelazador que fija la normalización se construye en el subespacio de
canales colectivos. Sean \(u_{90},u_{120}\) los vectores unitarios uniformes
de las incidencias hexádica y octádica, y sean \(v_{90},v_{120}\) los
vectores unitarios uniformes de los dos conjuntos de enlaces de
\eqref{p12-83-c14-s2:eq:corte-90-120-rev11}. El mapa
\begin{equation}
 \begin{aligned}
 \mathcal J_{6\to8}:{}
 &\operatorname{span}\{u_{90},u_{120}\}
 \longrightarrow
 \operatorname{span}\{v_{90},v_{120}\},\\
 &\mathcal J_{6\to8}u_{90}=v_{90},
 \qquad
 \mathcal J_{6\to8}u_{120}=v_{120}
 \end{aligned}
 \label{p12-83-c14-s2:eq:entrelazador-area-barbero-rev11}
\end{equation}
es una isometría entre estos subespacios colectivos. Si
\[
 D_{\rm inc}=90|u_{90}\rangle\langle u_{90}|
 +120|u_{120}\rangle\langle u_{120}|,
 \qquad
 D_{\rm area}=90|v_{90}\rangle\langle v_{90}|
 +120|v_{120}\rangle\langle v_{120}|,
\]
entonces
\[
 \boxed{\mathcal J_{6\to8}D_{\rm inc}
 =D_{\rm area}\mathcal J_{6\to8}.}
\]
En particular, la combinación primitiva \(12:-1\) actúa
con los mismos coeficientes antes y después de \(\mathcal J_{6\to8}\). Esta
igualdad fija el diccionario de normalización que inserta
\(\gamma_{\rm HMT}=\Gamma_{6\to8}\) en la conexión de
Ashtekar--Barbero.

La información retenida por la reducción a la región se publica mediante el
Hamiltoniano modular. Como estructura de partida de esta lectura se declara
\[
 \Delta_4^{\rm mod}
 =\frac{(A-C^*)\pi/180}{9\cdot90}
 \left(1-\frac7{12}\frac\pi{729}\right),
 \qquad w_{90}=90\Delta_4^{\rm mod},\quad
 w_{120}=120\Delta_4^{\rm mod},
\]
se define
\[
 K_{\rm link}(w)=wN,
 \qquad
 \rho_{\rm link}(w)=
 \frac{\operatorname{diag}(1,e^{-w},e^{-2w})}
 {1+e^{-w}+e^{-2w}},
\]
\begin{equation}
 \rho_A=
 \bigotimes_{e\in\partial A_{90}}\rho_{\rm link}(w_{90})
 \otimes
 \bigotimes_{e\in\partial A_{120}}\rho_{\rm link}(w_{120}),
 \qquad K_A=-\log\rho_A.
 \label{p12-83-c14-s2:eq:estado-hamiltoniano-modular-barbero-rev11}
\end{equation}
Entonces
\[
 K_A=\sum_{e\in\partial A}w_eN_e+c\mathbf1,
 \qquad
 S(\rho_A)=18S(w_{90})+18S(w_{120}),
\]
y la evaluación de esa estructura de partida da
\[
 S(\rho_A)=39.54837320881808\ldots\ \text{nats}
 =57.05624190358778\ldots\ \text{bits}.
\]
La cantidad \(S(\rho_A)\) es la entropía de von Neumann del producto mixto
declarado. Una purificación global la identifica con una entropía de
entrelazamiento a través del corte; sin esa purificación, su tipo canónico es
entropía modular reducida. Así, \(\gamma_{\rm HMT}\) fija el cuanto de
área/torsión, \(S_{\rm tri}\) cuenta la capacidad areal, \(S_{\rm bi}\)
cuantifica el reparto entre hojas y \(S(\rho_A)\) mide la mezcla modular.
Operador de área, espectro y entrelazador pertenecen a una sola construcción
de borde.
